%% kapa: exam=1 round=28 year=2017 period=1 form=A subject=real_estate subject_ko=부동산학원론 seq=15 number=95
%% source: Q-Net 감정평가사 기출문제 (cst00309) artlSeq=5199603 file=extracted/1008/5199603/2017년 제28회 감정평가사 1차시험 문제지/1교시 A형.pdf
%% conversion: kordoc 4.24.51 -> tools/gen1.py (auto); review=proofread-sonnet; flags=none
\begin{question}{95}
부동산의 개념 등에 관한 설명으로 옳지 않은 것은?
\begin{choices}
  \item 부동산이란 토지 및 그 정착물을 말하며, 부동산 이외의 물건은 동산이다.
  \item 부동산의 복합개념은 부동산을 법률적ㆍ경제적ㆍ기술적인 측면 등으로 이해하고자 하는 것이다.
  \item 부동산은 20년간 소유의 의사로 평온, 공연하게 점유하고 등기함으로써 그 소유권을 취득한다.
  \item 동산은 용익물권과 담보물권의 설정이 가능하다.
  \item 넓은 의미의 부동산에는 등기ㆍ등록의 대상이 되는 항공기ㆍ선박ㆍ자동차 등도 포함된다.
\end{choices}
\end{question}
\answer{4}
